\begin{solapartat}
La força gravitatòria és conservativa, per tant l'energia total d'un cos es conserva al llarg de la seva trajectòria: \punts{0,2}
% A l'original, la massa m del meteorit apareix ratllada (simplificada) a la primera equació.
\begin{align*}
&\frac{1}{2}\not{m}\,v_0^2 - G\,\frac{\not{m}M_L}{h_0 + R_L} = \frac{1}{2}\not{m}\,v_1^2 - G\,\frac{\not{m}M_l}{R_L} \ \text{\punts{0,3}}\\
&\Rightarrow \frac{1}{2}\,v_1^2 = \frac{1}{2}\,v_0^2 + G\,M_L\left\{\frac{1}{R_L} - \frac{1}{R_L + h_0}\right\} \Rightarrow
\end{align*}
\[ v_1 = \sqrt{v_0^2 + 2GM_L\,\frac{h_0}{R_L(R_L + h_0)}} \ \text{\punts{0,3}} = 4,71\cdot 10^{3}\ \mathrm{m/s} \ \text{\punts{0,2}} \]
\end{solapartat}

\begin{solapartat}
L'energia mecànica del meteorit serà:
\[ E_m = \frac{1}{2}\,m\,v_0^2 - G\,\frac{mM_L}{R_L + h_0} \ \text{\punts{0,2}} = 3,31\cdot 10^{9}\ \mathrm{J} \ \text{\punts{0,2}} \]
Per un cos de la mateixa massa, però en òrbita a la mateixa distància, l'energia mecànica és:
\[ E_o = -\frac{1}{2}\,G\,\frac{mM_L}{R_L + h_0} \ \text{\punts{0,2}} = -8,35\cdot 10^{7}\ \mathrm{J} \ \text{\punts{0,2}} \]
Com es pot comprovar: $E_m > E_o$ \punts{0,2}
\end{solapartat}
